\documentclass[10pt]{article}
\usepackage[margin=20mm]{geometry}
\usepackage{amsmath,amssymb,mathtools,fontspec,unicode-math,xurl,hyperref}
\IfFontExistsTF{DejaVu Serif}{\setmainfont{DejaVu Serif}}{\setmainfont{Latin Modern Roman}}
\setmathfont{latinmodern-math.otf}
\setlength{\parindent}{0pt}\setlength{\parskip}{6pt}\sloppy
\title{Technical Review of
Simple matrix expressions for the curvatures of Grassmannian}
\author{AI-assisted initial technical review}
\date{}
\begin{document}
\maketitle
Abstract: The involution representation gives useful explicit tangent, normal, second fundamental form and Riemann-curvature formulas. Direct checks nevertheless establish three Major issues: the claimed local product of many geodesic two-spheres fails a curvature-operator rank obstruction; the third fundamental form omits the dimension multiplying averaged mean curvature; and the Weyl convention is incompatible with the stated Riemann sign, producing an incorrect nonzero Bach tensor. Four Minor groups concern scalar/delta normalization, a Jacobi factor, tensor and bundle definitions, and a quotient-model comparison. The report supplies exact examples and distinguishes valid foundational formulas from consequences requiring revision.\par
\section*{Identity and review scope}
The reviewed source is Zehua Lai, Lek-Heng Lim and Ke Ye, Simple matrix expressions for the curvatures of Grassmannian, Foundations of Computational Mathematics 26 (2026), 2053–2089, published online 27 August 2025, DOI 10.1007/s10208-025-09723-9 (https://doi.org/10.1007/s10208-025-09723-9) [paper]. The assigned identifier is M25-0467. The acquired file is the published-format 37-page PDF hosted on the author’s institutional site; the online year and the later volume year are distinct metadata facts.\par
All 37 pages, including all definitions, proofs, model comparisons and 54 bibliography entries, were read. Actual page images inspected were physical pages 1, 3–4, 8–15, 17–30, 32–35. Physical page \(p\) corresponds to printed page \(2052+p\); locations below use the printed pagination. No separate code or supplement was confirmed in the supplied bounded acquisition record, and their absence is not treated as a mathematical defect. Reviewer-written exact matrix calculations check the examples below; no author code was executed. External cited theorems and priority claims were not all independently verified.\par
\section*{What the basic matrix calculation establishes}
Write \(l=n-k\), \(m=kl\), and work at \(Q=\operatorname{diag}(I_k,-I_l)\). For \(B\in\mathbb R^{k\times l}\), let \[X(B)=\begin{pmatrix}0&B\\B^T&0\end{pmatrix}.\] Differentiating \(Q^2=I\) gives this tangent space. Its Frobenius orthogonal complement consists of block-diagonal symmetric matrices. The metric is \(g(X(B),X(C))=2\mathop{\mathrm{tr}}(B^TC)\). The second fundamental form in Theorem 4 is \[\mathrm{II}(X(B),X(C))
 =\tfrac12\operatorname{diag}(-BC^T-CB^T,\ B^TC+C^TB).\] It is consistent with differentiating the defining equation twice. It immediately gives zero relative nullity, since \(\mathrm{II}(X,X)=0\) forces \(B=0\). The shape operator along \(H=\operatorname{diag}(H_1,H_2)\) acts by \[B\longmapsto\tfrac12(BH_2-H_1B),\] so its eigenvalues are \((\lambda_{k+j}-\lambda_i)/2\). The mean vector in Corollary 7 is the average \(m^{-1}\mathop{\mathrm{tr}}\mathrm{II}\), with the stated block entries. These results support the principal and directional Gaussian curvature formulas.\par
The Gauss equation, with the manuscript’s Riemann convention, gives \[R(X,Y,Z,W)=\tfrac12\mathop{\mathrm{tr}}([X,Y]ZW),\qquad
 R(X,Y,Y,X)=\tfrac14\|[X,Y]\|_F^2.\] Contraction in the orthonormal basis \(X(E_{ij}/\sqrt2)\) gives \[\mathop{\mathrm{Ric}}=\lambda g,\qquad \lambda=\frac{n-2}{8},\qquad
 \mathop{\mathrm{tr}}_g\mathop{\mathrm{Ric}}=m\lambda.\] These formulas pass the exact checks. They also furnish internal tests of several later expressions, without needing a different metric normalization or a different Grassmannian model.\par
\section*{Major issues}
Problem 1 (Major: the local sphere-product theorem is false in an allowed case).  Theorem 23(b), pp. 2080–2081, permits \(k=2,n=6,r=4\). It therefore claims an isometric totally geodesic embedding of an open subset of a product of four round two-spheres into \(\mathop{\mathrm{Gr}}(2,6)\). Both dimensions are eight. Under the paper’s convention that Riemannian embeddings are isometric, such an embedding would be a local isometry.\par
This is impossible by the curvature operator. On \(\Lambda^2T_Q\mathop{\mathrm{Gr}}(2,6)\), the curvature bilinear form is, up to the harmless ordering sign convention, the positive Gram form \[(X\wedge Y,Z\wedge W)\longmapsto
          \tfrac14\langle[X,Y],[Z,W]\rangle_F.\] The commutators span \(\mathfrak{so}(2)\oplus\mathfrak{so}(4)\), of dimension \(1+6=7\). Indeed, tangent basis pairs with a common column generate the first block, and pairs with a common row generate each skew elementary matrix in the second block. Consequently the curvature operator has rank seven. Exact evaluation of the \(28\times28\) Gram matrix confirms this rank.\par
For a product of four round two-spheres, each factor contributes one nonzero curvature-operator direction and every mixed two-plane has zero curvature. Its curvature operator has rank four, independently of the positive radii of the factors. A local isometry cannot change this rank. Thus the issue is a counterexample to the stated Riemannian conclusion, not just an omitted proof of the exponential-chart construction.\par
The proposed map \(\psi\circ\rho\circ\phi^{-1}\) is locally injective after shrinking a chart, but injectivity and orthogonality of selected tangent planes do not prove that it preserves the product metric. Mixed curvature and the relevant Lie-triple-system conditions are missing. In particular, the discussion contrasting the oriented and unoriented Grassmannians cannot remove a local curvature obstruction: a covering has the same local geometry.\par
Part (a) also needs a range correction unless an assumption \(k\le n/2\) is imposed. For \(k=3,n=4,r=1\), its printed bound holds, but its construction requires a one-dimensional \(W_0\) in the zero-dimensional complement of a four-dimensional block. Moreover \(\mathop{\mathrm{Gr}}(3,4)\) is real projective three-space with a scaled round metric. A complete totally geodesic surface lifts to a great two-sphere in its round universal cover, whose projected image is a projective plane; it cannot be an embedded two-sphere.\par
The valid orthogonal-block construction requires \[r\le \left\lfloor\frac{\min(k,n-k)}2\right\rfloor.\] This ensures that the fixed \((k-2r)\)-plane fits in the remaining \(n-4r\) dimensions. It supplies a conservative corrected global construction using a geodesic sphere in each \(\mathop{\mathrm{Gr}}(2,4)\) block. It does not establish the much larger local bound asserted in part (b).\par
Problem 2 (Major: the third fundamental form uses averaged mean curvature without its dimension).  The definition on p. 2061 is \(\mathrm{III}(X,Y)=\sum_\alpha g(S_{\eta_\alpha}X,S_{\eta_\alpha}Y)\), where the normal basis is orthonormal. In particular, \(\mathrm{III}(X,X)\ge0\). But the proof of Corollary 9, pp. 2071–2072, uses \(\mathrm{III}=\langle\mathrm{II},H\rangle-\mathop{\mathrm{Ric}}\), although Corollary 7 defines \(H=m^{-1}\mathop{\mathrm{tr}}\mathrm{II}\).\par
Contracting the Gauss equation directly gives \[\mathop{\mathrm{Ric}}(X,Y)=\langle\mathrm{II}(X,Y),\mathop{\mathrm{tr}}\mathrm{II}\rangle-\mathrm{III}(X,Y),\] so the required formula is \[\mathrm{III}(X,Y)=m\langle\mathrm{II}(X,Y),H\rangle-\mathop{\mathrm{Ric}}(X,Y)
          =\frac{n+2}{8}g(X,Y)
          =\frac{n+2}{4}\mathop{\mathrm{tr}}(B^TC).\] For \(\mathop{\mathrm{Gr}}(2,4)\) and unit \(X=X(E_{11}/\sqrt2)\), summing the six normal-basis shape-operator squares gives \(3/4\). Corollary 9 and Table 1 give zero. For some larger dimensions their coefficient is even negative. This contradicts the paper’s own definition and changes one of its principal new extrinsic formulas. The second fundamental form, averaged mean vector and Ricci tensor need not be changed to obtain the correction.\par
Problem 3 (Major: the Weyl sign produces a tensor with nonzero Ricci trace).  Let \[K_g(X,Y,Z,W)=g(X,Z)g(Y,W)-g(X,W)g(Y,Z).\] For an orthonormal basis \(e_a\), \[\sum_aK_g(X,e_a,e_a,Y)=-(m-1)g(X,Y).\] The Riemann convention in the paper gives \(\sum_aR(X,e_a,e_a,Y)=\lambda g(X,Y)\). Thus the tensor in Corollary 17, \[W_{\rm printed}=R-\frac{\lambda}{m-1}K_g,\] has contraction \(2\lambda g\), rather than zero. Its algebra follows the sign-incompatible definition on p. 2064, but that definition does not yield the Weyl tensor for the previously fixed Riemann ordering.\par
The Einstein-case correction is \[W=R+\frac{\lambda}{m-1}K_g.\] The corresponding Ricci-decomposition signs must be changed consistently in Section 3 and the conclusion. As a concrete four-dimensional test, take \(\mathop{\mathrm{Gr}}(1,5)\). The matrix formula gives constant sectional curvature \(1/8\), \(\lambda=3/8\), and the corrected \(W\) vanishes in every component. The printed \(W\) has contraction \(3g/4\).\par
This error propagates to Corollary 18 and Table 1. In this example the genuine Bach tensor is zero, whereas the printed formula gives \[B_{\rm printed}=\frac9{64}g.\] More generally, with the genuine trace-free Weyl tensor and the parallel Einstein Ricci tensor, the Cotton/derivative contribution and the Ricci contraction in the Bach formula vanish, in dimensions where that formula is defined. Hence the claimed nonzero Bach expression must be replaced by zero. The Codazzi conclusion remains true for this corrected zero tensor; it does not validate the old value.\par
\section*{Further corrections}
Problem 4 (Minor: scalar and delta conventions, and the maximizing basis).  Page 2063 identifies \(\mathop{\mathrm{tr}}\mathop{\mathrm{Ric}}\) with the sum of sectional curvatures over unordered pairs. The correct identity is \[\mathop{\mathrm{tr}}\mathop{\mathrm{Ric}}=2\sum_{a<b}\kappa(e_a,e_b).\] The same page uses the unordered-pair sum to define the scalar curvature of a subspace, then asserts that the whole tangent space has the previously defined full scalar curvature. The two conventions cannot both be retained.\par
For \(m=k(n-k)\), using the unordered-pair scalar consistently gives the delta values \[\overline\delta_{2,r}=\frac{m(n-2)}{16},\qquad
 \underline\delta_{2,r}=\frac{m(n-2)}{16}-\frac r4\] in the intended range. Using full Ricci traces consistently for the ambient tangent space and its two-planes instead gives \[\overline\delta_{2,r}=\frac{m(n-2)}8,\qquad
 \underline\delta_{2,r}=\frac{m(n-2)}8-\frac r2.\] The printed pair mixes the full ambient trace with half-traces on the selected planes. It can describe a specifically named hybrid quantity, but not both scalar conventions at once. For \(\mathop{\mathrm{Gr}}(2,4)\), the unordered scalar is \(1/2\), the full scalar is 1, and the lower invariant at \(r=2\) is zero under either consistent convention; the printed expression is \(1/2\).\par
There is also a repairable error in the maximizing basis (16), p. 2076. In block notation its two pairs are \[(E_{11}/\sqrt2,\ J/2),\qquad(E_{22}/\sqrt2,\ K/2),
 \quad
 J=\begin{pmatrix}0&1\\1&0\end{pmatrix},\
 K=\begin{pmatrix}0&1\\-1&0\end{pmatrix}.\] Their sectional curvatures are \(1/8\) each, not \(1/4\). Use instead \[(I_2/2,\ K/2),\qquad(D/2,\ J/2),
 \quad D=\operatorname{diag}(1,-1).\] The four full tangent matrices are orthonormal, and each pair has curvature \(1/4\). These blocks repair the attainment calculation for the sectional sums. The original zero-curvature pairs pass the direct check. This is a normalization and witness correction, not a rejection of the block extremization idea.\par
Problem 5 (Minor: the Jacobi expression is twice its defining symmetrization).  The definition on p. 2062 is \(J(X,Y,Z,W)=\tfrac12(R(X,Y,Z,W)+R(Z,Y,X,W))\). Substitution of Proposition 10 yields \[J(X,Y,Z,W)
 =\tfrac12\mathop{\mathrm{tr}}(XYZW)
       -\tfrac14\mathop{\mathrm{tr}}\bigl(Y(XZ+ZX)W\bigr).\] Corollary 11, p. 2073, and Table 1 give twice this expression. On \(\mathop{\mathrm{Gr}}(1,3)\), take the orthonormal tangent matrices corresponding to \(B=(1,0)/\sqrt2\) and \(C=(0,1)/\sqrt2\). Then the definition gives \(J(X,Y,X,Y)=-1/8\), while the displayed formula gives \(-1/4\). Multiply both its full-matrix and block formulas by \(1/2\).\par
Problem 6 (Minor: tensor divergence, bundle language and dimension restrictions).  The rule on p. 2062 that differentiates only the final factor of a decomposable tensor is not well-defined over \(C^\infty(M)\). On the real line, with \(e=\partial_x\), \[(xe)\otimes e=e\otimes(xe).\] The printed rule assigns zero to the first expression and \(e\) to the second. Tensor divergence should contract the covariant derivative of the complete tensor, including derivatives of all factors.\par
In Proposition 1, p. 2065, vertical cocurvature vanishes because the vertical distribution is tangent to the fibers and hence integrable. The proof instead calls the horizontal distribution integrable with integral manifold \(M\); for a connection with nonzero curvature this is false. Use horizontal lifts and the pullback tangent bundle, rather than a literal global inclusion \(\mathfrak X(M)\subset\mathfrak X(O(M))\). The stated vanishing of torsion, nonmetricity and vertical cocurvature can be retained after this repair.\par
Finally, formulas containing \(m-2\) or \(m-3\) must carry dimension restrictions. Use the standard Schouten/Cotton definitions for \(m>2\) and the displayed Bach definition for \(m>3\), or state separate extensions. Cancellation of a denominator in a special Einstein expression does not define the original general tensor in excluded dimensions.\par
Problem 7 (Minor: the general-linear quotient and common-point lifts).  Equation (17), p. 2082, includes \(\mathrm{GL}(n)/(\mathrm{GL}(k)\times\mathrm{GL}(n-k))\). Its dimension is \(2k(n-k)\), twice the dimension of the Grassmannian. The stabilizer of a \(k\)-plane is the block upper-triangular parabolic subgroup, including the off-diagonal \(k\times(n-k)\) block. Substituting that stabilizer repairs the quotient identification.\par
The subsequent assertion that horizontal lifts must emerge at different representatives is also not inherent. Once \(Q\) is fixed, the map \((d_Q\pi)^{-1}\) used in the paper sends every tangent vector to the horizontal subspace of that same \(T_QO(n)\). Four vectors can therefore be lifted at one chosen \(Q\). Arbitrarily supplied incompatible representatives may need alignment, but a common representative does not.\par
These corrections leave room for practical advantages of the involution model. Such advantages should be tied to specified operations or algorithms. The affine relation \(Q=2P-I\) gives \(dQ=2\,dP\) and a factor of four in the induced metric; the singularity of a projector as a square matrix alone is not an analysis of the numerical stability of every projector-based algorithm. No comparative numerical experiment was performed in this review.\par
\section*{Assessment}
Major revision is required by the established sphere-product obstruction and the incorrect geometric tensors, independently of the smaller corrections. The verified tangent, normal, second-form, shape and basic curvature calculations provide a substantial usable core. The report therefore does not suggest that every formula in Table 1 fails. Corrections should be propagated consistently through that table, definitions, proofs and geometric consequences, with the Riemann convention and the full-versus-half scalar convention fixed before contraction.\par
The exact checks include the complete curvature-operator Gram matrix, shape-square sums, Weyl contractions, Jacobi comparison and both the printed and repaired extremizing pairs. They are targeted checks, not a formal verification of every theorem or a reproduction of a numerical experiment. Scientific adjudication remains pending. This is an initial technical review and has not undergone independent acceptance.\par
1 Z. Lai, L.-H. Lim and K. Ye, Simple matrix expressions for the curvatures of Grassmannian, Found. Comput. Math. 26 (2026), 2053–2089; online 27 August 2025. https://doi.org/10.1007/s10208-025-09723-9. Acquired published-format PDF: https://www.stat.uchicago.edu/\textasciitilde{}lekheng/work/curv.pdf.\par
\end{document}
